\documentclass[../../main.tex]{subfiles}% 注意这里的文件路径不能用 ./main.tex ,否则用latexmk编译子文件会报错
\graphicspath{{\subfix{./image/}}{\subfix{../../image/}}} % 指定图片引用路径,后续可以直接使用图片文件名
% 如果图片存放在上级目录,则使用 ../../image/ 作为备用路径

% 例如:
% \begin{figure}[H]
% \centering
% \includegraphics[width=0.3\textwidth]{图.png}
% \caption{}
% \label{figure:图}
% \end{figure}
% 注意:上述\label{}一定要放在\caption{}之后,否则引用图片序号会只会显示??.

\begin{document}

\section{有理标准型的几何与应用}

回顾\hyperref[section:有理标准型]{有理标准型和循环子空间相关理论}.

\begin{proposition}\label{proposition:矩阵有不同特征值则特征多项式和极小多项式一定相等}
设 $n$ 阶矩阵 $A$ 有 $n$ 个不同的特征值，求证：$A$ 的特征多项式和极小多项式相等。
\end{proposition}
\begin{proof}

{\color{blue}证法一:}
设 $A$ 的 $n$ 个不同的特征值为 $\lambda_1,\lambda_2,\cdots,\lambda_n$，则由\refcor{corollary:极小多项式和特征多项式有相同的根(不计重数)}可知，特征多项式 $f(\lambda)$ 和极小多项式 $m(\lambda)$ 有相同的根（不计重数），因此 
\begin{align*}
f(\lambda)=m(\lambda)=(\lambda - \lambda_1)(\lambda - \lambda_2)\cdots(\lambda - \lambda_n)
\end{align*}

{\color{blue}证法二:}
由于 $A$ 有 $n$ 个不同的特征值，故 $A$ 相似于对角矩阵。又因为相似矩阵有相同的特征多项式和极小多项式，所以只要对对角矩阵证明此结论即可。设 $A = \mathrm{diag}\{\lambda_1,\lambda_2,\cdots,\lambda_n\}$，则 $\lambda I_n - A = \mathrm{diag}\{\lambda - \lambda_1,\lambda - \lambda_2,\cdots,\lambda - \lambda_n\}$，这是一个主对角元素两两互素的对角矩阵，由\hyperref[theorem:lambda-矩阵和初等因子的基本性质]{$\lambda$-矩阵和初等因子的基本性质(1)}以及数学归纳法可知其法式为 $\mathrm{diag}\{1,\cdots,1,(\lambda - \lambda_1)(\lambda - \lambda_2)\cdots(\lambda - \lambda_n)\}$。因此，$A$ 的特征多项式和极小多项式相等.
\end{proof}

\begin{proposition}\label{proposition:不同特征值对应的特征向量之和恰好构成循环向量}
设$n$阶矩阵$A$有$n$个不同的特征值,且特征值 $\lambda_i$ 对应的特征向量为 $\alpha_i$，由\refcor{corollary:线性变换属于不同特征值的特征向量必线性无关}可知$\{\alpha_1,\cdots,\alpha_n\}$ 为 $\mathbb{C}^n$ 的一组基。则$\alpha = \alpha_1 + \cdots + \alpha_n$ 是 $A$ 的循环空间 $\mathbb{C}^n$ 的循环向量,即$\mathbb{C}^n = L(\alpha,A\alpha,\cdots,A^{n - 1}\alpha)=C(A,\alpha)$ 为循环空间，$\alpha$ 是循环向量。
\end{proposition}
\begin{proof}
事实上，由 $A^k\alpha = \lambda_1^k\alpha_1 + \cdots + \lambda_n^k\alpha_n$，利用 Vandermonde 行列式容易证明 $\{\alpha,A\alpha,\cdots,A^{n - 1}\alpha\}$ 是 $\mathbb{C}^n$ 的一组基，从而 $\mathbb{C}^n = L(\alpha,A\alpha,\cdots,A^{n - 1}\alpha)=C(A,\alpha)$ 为循环空间，$\alpha$ 是循环向量。
\end{proof}

\begin{proposition}\label{proposition:特征多项式在数域K上的不可约分解则在复数域上可对角化}
设数域 $\mathbb{K}$ 上的 $n$ 阶矩阵 $A$ 的特征多项式 $f(\lambda)=P_1(\lambda)P_2(\lambda)\cdots P_k(\lambda)$，其中 $P_i(\lambda) (1\leqslant  i\leqslant  k)$ 是 $\mathbb{K}$ 上互异的首一不可约多项式。求证：$A$ 的有理标准型只有一个 Frobenius 块，并且 $A$ 在复数域上可对角化。
\end{proposition}
\begin{remark}
我们也可以利用\refthe{theorem:F(A)=C(A)的充要条件和不变子空间刻画}和初等因子证明这个命题。若利用不变因子在基域扩张下的不变性，则这个命题也可由\refpro{proposition:矩阵有不同特征值则特征多项式和极小多项式一定相等}得到。
\end{remark}
\begin{proof}
设 $A$ 的不变因子组为 $d_1(\lambda),d_2(\lambda),\cdots,d_n(\lambda)$，其中$d_{i}(\lambda)\mid d_{i+1}(\lambda),i=1,2,\cdots,n-1$,则有
\begin{align}\label{equation--proposition0.3--1.1.0}
f(\lambda)=P_1(\lambda)P_2(\lambda)\cdots P_k(\lambda)=d_1(\lambda)d_2(\lambda)\cdots d_n(\lambda)
\end{align}
由于 $P_i(\lambda)$ 是不可约多项式，故存在某个 $j$，使得 $P_i(\lambda)\mid d_j(\lambda)$，否则,由\hyperref[proposition:不可约多项式的基本性质]{不可约多项式的基本性质(1)}可知$(P_i(\lambda)$,$d_j$ $(\lambda))$ $=1$,$j=1,2,\cdots,n$.再由\hyperref[proposition:互素多项式和最大公因式的基本性质]{互素多项式和最大公因式的基本性质(5)}可知$(P_i(\lambda),d_1(\lambda)d_2(\lambda)\cdots d_n(\lambda))=1$,这与\eqref{equation--proposition0.3--1.1.0}矛盾!
从而 $P_i(\lambda)\mid d_n(\lambda) (1\leqslant  i\leqslant  k)$。由\hyperref[proposition:互素多项式和最大公因式的基本性质]{互素多项式和最大公因式的基本性质(1)}可知，$P_1(\lambda)P_2(\lambda)\cdots P_k(\lambda)\mid d_n(\lambda)$，因此只能是 $d_1(\lambda)=\cdots = d_{n - 1}(\lambda)=1$，$d_n(\lambda)=f(\lambda)$，从而 $A$ 的有理标准型只有一个 Frobenius 块。由于特征多项式 $f(\lambda)=P_1(\lambda)P_2(\lambda)\cdots P_k(\lambda)$ 在 $\mathbb{K}$ 上无重因式，故 $(f(\lambda),f'(\lambda)) = 1$，从而 $f(\lambda)$ 在复数域上无重根，即 $A$ 有 $n$ 个不同的特征值，于是 $A$ 在复数域上可对角化。
\end{proof}

\begin{corollary}
设数域 $\mathbb{K}$ 上的 $n$ 阶矩阵 $A$ 的特征多项式 $f(\lambda)=P_1(\lambda)P_2(\lambda)\cdots P_k(\lambda)$，其中 $P_i(\lambda) (1\leqslant  i\leqslant  k)$ 是 $\mathbb{K}$ 上互异的首一不可约多项式.并且 $\alpha_i$ 为线性方程组 $P_i(A)x = 0$ 的非零解，则 $\alpha=\alpha_1 + \cdots + \alpha_k$ 是 $A$ 的循环空间 $\mathbb{K}^n$ 的循环向量。
\end{corollary}
\begin{proof}
由\refpro{proposition:特征多项式在数域K上的不可约分解则在复数域上可对角化}及\refthe{theorem:不变子空间是循环子空间的充要条件}可知$\mathbb{K}^n$就是一个循环空间.(未完成证明)
\end{proof}

\begin{proposition}\label{proposition:特征多项式是不可约多项式的刻画}
设 $\varphi$ 是数域 $\mathbb{K}$ 上 $n$ 维线性空间 $V$ 上的线性变换，$\varphi$ 的特征多项式为 $f(\lambda)$，证明以下 3 个结论等价：
\begin{enumerate}[(1)]
\item $V$ 只有平凡的 $\varphi$-不变子空间；

\item $V$ 中任一非零向量都是循环向量，使 $V$ 成为循环空间；

\item $f(\lambda)$ 是 $\mathbb{K}$ 上的不可约多项式。
\end{enumerate}
\end{proposition}
\begin{proof}
$(1)\Rightarrow(2)$：任取 $V$ 中非零向量 $\alpha$，则循环子空间 $C(\varphi,\alpha)$ 是非零 $\varphi$-不变子空间。由于 $V$ 只有平凡的 $\varphi$-不变子空间，故 $C(\varphi,\alpha)=V$，即 $V$ 中任一非零向量都是循环向量，使 $V$ 成为循环空间。

$(2)\Rightarrow(3)$：用反证法，假设 $f(\lambda)=g(\lambda)h(\lambda)$，其中 $g(\lambda),h(\lambda)$ 是 $\mathbb{K}$ 上次数小于 $n$ 的首一多项式。由 Cayley - Hamilton 定理可知 $0=f(\varphi)=g(\varphi)h(\varphi)$，取行列式可知$g(\varphi),h(\varphi)$ 中至少有一个是奇异线性变换，不妨设为 $g(\varphi)$，则 $\mathrm{Ker}g(\varphi)\neq \{0\}$.任取 $\mathrm{Ker}g(\varphi)$ 中的非零向量 $\alpha$，设 $\deg g(\lambda)=r$，则不妨设
\begin{align*}
g(\varphi) = a_r \varphi^r + a_{r-1} \varphi^{r-1} + \cdots + a_1, \quad \text{其中} \ a_r \neq 0.
\end{align*}
由$\alpha \in \mathrm{Ker} g(\varphi)$可知
\begin{align*}
g(\varphi)(\alpha) = a_r \varphi^r(\alpha) + a_{r-1} \varphi^{r-1}(\alpha) + \cdots + a_1 \alpha = 0.
\end{align*}
于是
\begin{align}
\varphi^r(\alpha) = -\frac{a_{r-1}}{a_r} \varphi^{r-1}(\alpha) - \cdots - \frac{a_1}{a_r} \alpha. \label{equation--:_:_:_:_::0.1}
\end{align}
假设对$k \geqslant  r$且$k \in \mathbb{N}$，成立$\varphi^k(\alpha)$可由$\{ \alpha, \varphi(\alpha), \cdots, \varphi^{r-1}(\alpha) \}$线性表示，则对\eqref{equation--:_:_:_:_::0.1}式两边同时作用$\varphi^{k-r+1}$可得
\begin{align*}
\varphi^{k+1}(\alpha) = -\frac{a_{r-1}}{a_r} \varphi^k(\alpha) - \cdots - \frac{a_1}{a_r} \varphi^{k-r+1}(\alpha).
\end{align*}
于是由归纳假设可知，$\varphi^{k+1}(\alpha)$可由$\{ \alpha, \varphi(\alpha), \cdots, \varphi^{r-1}(\alpha) \}$线性表示. 故由数学归纳法可得，对$\forall k \geqslant  r$且$k \in \mathbb{N}$，都有$\varphi^k(\alpha)$可由$\{ \alpha, \varphi(\alpha), \cdots, \varphi^{r-1}(\alpha) \}$线性表示.因此
$C(\varphi ,\alpha )=L(\alpha ,\varphi (\alpha ),\cdots \cdots )=L(\alpha ,\varphi (\alpha ),\cdots ,\varphi ^{r-1}(\alpha ))$，其维数 $\leqslant  r < n$，故 $C(\varphi,\alpha)\neq V$，这与 $V$ 中任一非零向量都是循环向量矛盾!

$(3)\Rightarrow(1)$：用反证法，假设存在非平凡的 $\varphi$-不变子空间 $U$，$\dim U = r$，则 $\varphi$ 在一组基下的表示矩阵为分块上三角矩阵 $M = \begin{pmatrix}
\boldsymbol{A} & \boldsymbol{C} \\
\boldsymbol{O} & \boldsymbol{B}
\end{pmatrix}$，其中 $\boldsymbol{A}$ 是 $\varphi|_U$ 的表示矩阵。于是特征多项式
\begin{align*}
f(\lambda)=|\lambda\boldsymbol{I}_V - \varphi|=|\lambda\boldsymbol{I}_n - M|=|\lambda\boldsymbol{I}_r - \boldsymbol{A}|\cdot|\lambda\boldsymbol{I}_{n - r} - \boldsymbol{B}|.
\end{align*}
是两个低次多项式的乘积，这与 $f(\lambda)$ 的不可约性矛盾！
\end{proof}

\begin{proposition}\label{proposition:极小多项式是不可约多项式的刻画}
设 $\varphi$ 是数域 $\mathbb{K}$ 上 $n$ 维线性空间 $V$ 上的线性变换，$\varphi$ 的极小多项式为 $m(\lambda)$。证明：$m(\lambda)$ 是 $\mathbb{K}$ 上的不可约多项式的充要条件是 $V$ 的任一非零 $\varphi$-不变子空间 $U$ 必为如下形式：
\begin{align*}
U = C(\varphi,\alpha_1)\oplus C(\varphi,\alpha_2)\oplus\cdots\oplus C(\varphi,\alpha_k)
\end{align*}
并且 $\varphi|_{C(\varphi,\alpha_i)}$ 的极小多项式都是 $m(\lambda)$。此时，$\varphi|_U$ 的极小多项式也是 $m(\lambda)$。
\end{proposition}
\begin{proof}

{\heiti 必要性：}设 $\varphi|_U$ 的极小多项式为 $n(\lambda)$，则 $m(\varphi|_U)=m(\varphi)|_U = 0$，从而 $n(\lambda)\mid m(\lambda)$。因为 $m(\lambda)$ 不可约，所以 $n(\lambda)=m(\lambda)$。又由于 $\varphi|_U$ 的所有不变因子都要整除 $m(\lambda)$ 且 $m(\lambda)$ 不可约，故所有的非常数不变因子都等于 $m(\lambda)$。最后，由\hyperref[theorem:有理标准型的几何意义]{有理标准型的几何意义}即得
\begin{align*}
U = C(\varphi,\alpha_1)\oplus C(\varphi,\alpha_2)\oplus\cdots\oplus C(\varphi,\alpha_k).
\end{align*}
并且\(\varphi|_{C(\varphi,\alpha_i)}\) 在基 \(\{\alpha_i,\varphi(\alpha_i),\cdots,\varphi^{r_i - 1}(\alpha_i)\}\) 下的表示矩阵就是友阵 \(C(m(\lambda))\)，其中 \(r_i = \dim C(\varphi,\alpha_i)\)。于是\(\varphi|_{C(\varphi,\alpha_i)}\)的极小多项式就是其表示矩阵\(C(m(\lambda))\)的极小多项式.又由\reflem{lemma:Frobenius标准型矩阵的极小多项式和不变因子}可知,\(C(m(\lambda))\)的极小多项式就是$m(\lambda)$,并且$n(\lambda)=m(\lambda)$,
故结论得证.

{\heiti 充分性：}用反证法，设 $m(\lambda)=g(\lambda)h(\lambda)$，其中 $g(\lambda),h(\lambda)$ 是 $\mathbb{K}$ 上次数小于 $m(\lambda)$ 次数的首一多项式，则 $0=m(\varphi)=g(\varphi)h(\varphi)$，取行列式可知 $g(\varphi),h(\varphi)$ 中至少有一个是奇异线性变换，不妨设为 $g(\varphi)$，于是 $\mathrm{Ker}g(\varphi)\neq \{0\}$。任取 $\mathrm{Ker}g(\varphi)$ 中的非零向量 $\alpha$，得到循环子空间 $U = C(\varphi,\alpha)$，
由 $g(\varphi)(\alpha)=0$可知
\begin{align*}
g(\varphi )(\varphi ^k\left( \alpha \right) )=\varphi ^k\left( g\left( \varphi \right) \left( \alpha \right) \right) =0,\quad k=0,1,2,\cdots .
\end{align*}
因此$g(\varphi|_U)=g(\varphi)|_U = 0$，于是 $\varphi|_U$ 的极小多项式$m(\lambda)$整除 $g(\lambda)$，从而其次数 $\leqslant \deg g(\lambda)<\deg m(\lambda)$，这与条件矛盾！
\end{proof}

\begin{theorem}[基于初等因子组的有理标准型]\label{theorem:基于初等因子组的有理标准型}
设数域 $\mathbb{K}$ 上的 $n$ 阶矩阵 $A$ 的初等因子组为 $P_1(\lambda)^{r_1},P_2(\lambda)^{r_2},\cdots,P_k(\lambda)^{r_k}$，证明：$A$ 相似于分块对角矩阵
\begin{align*}
\widetilde{\boldsymbol{F}}&=\mathrm{diag}\{\boldsymbol{F}(P_1(\lambda)^{r_1}),\boldsymbol{F}(P_2(\lambda)^{r_2}),\cdots,\boldsymbol{F}(P_k(\lambda)^{r_k})\}\\
\widetilde{\boldsymbol{C}}&=\mathrm{diag}\{\boldsymbol{C}(P_1(\lambda)^{r_1}),\boldsymbol{C}(P_2(\lambda)^{r_2}),\cdots,\boldsymbol{C}(P_k(\lambda)^{r_k})\}
\end{align*}
称为 $A$ 的\textbf{基于初等因子组的有理标准型}。
\end{theorem}
\begin{proof}
由 Frobenius 块和友阵的性质可知，$\lambda\boldsymbol{I}_n - \widetilde{\boldsymbol{F}}$ 和 $\lambda\boldsymbol{I}_n - \widetilde{\boldsymbol{C}}$ 都相抵于
\begin{align*}
\mathrm{diag}\{1,\cdots,1,P_1(\lambda)^{r_1};1,\cdots,1,P_2(\lambda)^{r_2};\cdots;1,\cdots,1,P_k(\lambda)^{r_k}\}
\end{align*}
再由\hyperref[theorem:lambda-矩阵和初等因子的基本性质]{$\lambda$-矩阵和初等因子的基本性质(2)}可知，$\widetilde{\boldsymbol{F}},\widetilde{\boldsymbol{C}}$ 与 $A$ 有相同的初等因子组，从而它们相似。
\end{proof}

\begin{theorem}\label{theorem:基于初等因子组的循环子空间之和分解}
设 $\varphi$ 是数域 $\mathbb{K}$ 上 $n$ 维线性空间 $V$ 上的线性变换，$\varphi$ 的初等因子组为 $P_1(\lambda)^{r_1},P_2(\lambda)^{r_2},\cdots,P_k(\lambda)^{r_k}$。证明：存在 $\alpha_1,\alpha_2,\cdots,\alpha_k\in V$，使得
\begin{align*}
V = C(\varphi,\alpha_1)\oplus C(\varphi,\alpha_2)\oplus\cdots\oplus C(\varphi,\alpha_k)
\end{align*}
\end{theorem}
\begin{proof}
由\hyperref[theorem:基于初等因子组的有理标准型]{基于初等因子组的有理标准型}和\refthe{theorem:不变子空间是循环子空间的充要条件}即得。
\end{proof}

\begin{example}
求证：存在 $n$ 阶实方阵 $A$，满足 $A^2 + 2A + 5I_n = O$ 的充要条件是 $n$ 为偶数。
当 $n\geqslant  4$ 时，验证满足上述条件的矩阵 $A$ 有无限个不变子空间。
\end{example}
\begin{proof}

{\heiti 必要性：}注意到 $A$ 适合多项式 $g(\lambda)=\lambda^2 + 2\lambda + 5$，故 $A$ 的极小多项式 $m(\lambda)\mid g(\lambda)$，又因为 $g(\lambda)$ 在实数域上不可约，故只能是 $m(\lambda)=g(\lambda)$。同理可证 $A$ 所有的非常数不变因子都等于 $g(\lambda)$，从而 $A$ 的不变因子组为 $1,\cdots,1,g(\lambda),\cdots,g(\lambda)$（$k$ 个 $g(\lambda)$）。因此 $A$ 的特征多项式 $f(\lambda)=g(\lambda)^k$，于是 $n = \deg f(\lambda)=2k$ 为偶数。

{\heiti 充分性：}设 $n = 2k$ 为偶数，则由必要性的证明可知，$A$ 的不变因子组为 $1,\cdots,1,g(\lambda),\cdots,g(\lambda)$（$k$ 个 $g(\lambda)$）。可用有理标准型构造满足条件的矩阵：
\begin{align}\label{equation----Akkk}
A=\mathrm{diag}\left\{\begin{pmatrix}
0 & -5 \\
1 & -2
\end{pmatrix},\cdots,\begin{pmatrix}
0 & -5 \\
1 & -2
\end{pmatrix}\right\} \text{（$k$ 个二阶方阵）}
\end{align}
当 $n\geqslant  4$ 时，设 $\{e_1,e_2,e_3,e_4\}$ 是前 4 个标准单位列向量，则容易验证\eqref{equation----Akkk}式中的矩阵$A$满足$Ae_1=e_2,Ae_3=e_4$,于是构造循环子空间 $\{C_l := C(A,e_1 + le_3)=L(e_1 + le_3,e_2 + le_4),l\in\mathbb{R}\}$ ,进一步容易验证循环子空间 $\{C_l := C(A,e_1 + le_3)=L(e_1 + le_3,e_2 + le_4),l\in\mathbb{R}\}$ 是两两互异的 $A$-不变子空间，故 $A$ 有无限个不变子空间。
\end{proof}

\begin{proposition}\label{proposition:极小多项式的次数小于等于矩阵的秩加1}
设 $A$ 是数域 $\mathbb{K}$ 上的 $n$ 阶方阵，求证：$A$ 的极小多项式的次数小于等于 $\mathrm{r}(A)+1$。
\end{proposition}
\begin{proof}

{\color{blue}证法一:}
设 $A$ 的不变因子组为 $1,\cdots,1,d_1(\lambda),\cdots,d_k(\lambda)$，则极小多项式 $m(\lambda)=d_k(\lambda)$,并且由\refthe{theorem:有理标准型核心定理}可知$A$ 相似于 $\boldsymbol{F}$ $=$ $\mathrm{diag}$ $\{$ $\boldsymbol{F}$ $(d_1(\lambda))$,$\cdots$ ,$\boldsymbol{F}$ $(d_k(\lambda))$ $\}$.设 $\deg d_k(\lambda)=r$，若 $d_k(0)\neq 0$，则由\hyperref[lemma:Frobenius标准型矩阵的极小多项式和不变因子]{Frobenius块的基本性质(1)}可知 $\boldsymbol{F}(d_k(\lambda))$ 非异；若 $d_k(0)=0$，则由\hyperref[lemma:Frobenius标准型矩阵的极小多项式和不变因子]{Frobenius块的基本性质(1)}可知 $\boldsymbol{F}(d_k(\lambda))$ 奇异且右上角的 $r - 1$ 阶子式非零，从而秩为 $r - 1$。因此，$\mathrm{r}(A)=\mathrm{r}(\boldsymbol{F})\geqslant \mathrm{r}(\boldsymbol{F}(d_k(\lambda)))\geqslant  r - 1=\deg d_k(\lambda)-1$.

{\color{blue}证法二:}
从 $A$ 的极小多项式 $m(\lambda)$ 分离出来的初等因子中，形如 $\lambda^r$ 的初等因子至多只有 1 个，对应于零特征值的 Jordan 块 $J_r(0)$,其秩为$r-1$，其余的初等因子对应于非零特征值的 Jordan块,其秩都为其阶数.因此 $\mathrm{r}(A)$ 大于等于这些 Jordan 块秩的和，后者等于 $\mathrm{deg}\,m(\lambda) - 1$ 或 $\mathrm{deg}\,m(\lambda)$.
\end{proof}

\begin{proposition}\label{proposition:特征值对应的特征矩阵的秩与不变因子组的关系}
设数域 $\mathbb{K}$ 上的 $n$ 阶矩阵 $A$ 的不变因子组是 $1,\cdots,1,d_1(\lambda),\cdots,d_k(\lambda)$，其中 $d_i(\lambda)$ 是非常数首一多项式，$d_i(\lambda)\mid d_{i + 1}(\lambda) (1\leqslant  i\leqslant  k - 1)$。求证：对 $A$ 的任一特征值 $\lambda_0$,
\begin{align*}
\mathrm{r}(\lambda_0\boldsymbol{I}_n - A)=n - \sum_{i = 1}^{k}\delta_{d_i(\lambda_0),0}
\end{align*}
其中记号 $\delta_{a,b}$ 表示：若 $a = b$，取值为 $1$；若 $a\neq b$，取值为 $0$。
\end{proposition}
\begin{proof}

{\color{blue}证法一:}
设 $\deg d_i(\lambda)=r_i$，则由\refthe{theorem:有理标准型核心定理}可知$A$ 相似于 $\boldsymbol{F}=\mathrm{diag}\{\boldsymbol{F}(d_1(\lambda)),\cdots,\boldsymbol{F}(d_k(\lambda))\}$，而相似矩阵有相同的特征多项式，故$\lambda_0I_n-A=\lambda_0I_n-F$.由\hyperref[lemma:Frobenius标准型矩阵的极小多项式和不变因子]{Frobenius块的基本性质(2)}可知$|\lambda_0\boldsymbol{I}_{r_i}-\boldsymbol{F}(d_i(\lambda))|=d_i(\lambda_0)$。若 $d_i(\lambda_0)\neq 0$，则 $\lambda_0\boldsymbol{I}_{r_i}-\boldsymbol{F}(d_i(\lambda))$ 非异；若 $d_i(\lambda_0)=0$，则 $\lambda_0\boldsymbol{I}_{r_i}-\boldsymbol{F}(d_i(\lambda))$ 奇异且右上角的 $r_i - 1$ 阶子式非零，从而秩为 $r_i - 1$。因此，
\begin{align*}
\mathrm{r}(\lambda_0\boldsymbol{I}_n - A)&=\mathrm{r}(\lambda_0\boldsymbol{I}_n - \boldsymbol{F})=\sum_{i = 1}^{k}\mathrm{r}(\lambda_0\boldsymbol{I}_{r_i}-\boldsymbol{F}(d_i(\lambda)))\\
&=\sum_{i = 1}^{k}(r_i - \delta_{d_i(\lambda_0),0})=n - \sum_{i = 1}^{k}\delta_{d_i(\lambda_0),0}
\end{align*}

{\color{blue}证法二:}
由\refthe{theorem:有理标准型核心定理}可知存在可逆 $\lambda$-矩阵 $\boldsymbol{P}(\lambda),\boldsymbol{Q}(\lambda)$，使得
\begin{align*}
\boldsymbol{P}(\lambda)(\lambda\boldsymbol{I}_n - A)\boldsymbol{Q}(\lambda)=\mathrm{diag}\{1,\cdots,1,d_1(\lambda),\cdots,d_k(\lambda)\}
\end{align*}
在上式中令 $\lambda = \lambda_0$，注意到 $\boldsymbol{P}(\lambda_0),\boldsymbol{Q}(\lambda_0)$ 是 $\mathbb{K}$ 上的可逆矩阵，故 $\lambda_0\boldsymbol{I}_n - A$ 相抵于 $\mathrm{diag}\{1,\cdots,1,d_1(\lambda_0),\cdots,d_k(\lambda_0)\}$，于是 $\mathrm{r}(\lambda_0\boldsymbol{I}_n - A)$ 等于 $n$ 减去等于零的 $d_i(\lambda_0)$ 的个数，从而结论得证。
\end{proof}

\begin{proposition}[迹0矩阵与交换子]\label{proposition:迹零阵必相似于主对角全为零的矩阵}
\begin{enumerate}[(1)]
\item\label{proposition:迹零阵必相似于主对角全为零的矩阵-1} 设 $A$ 是数域 $\mathbb{K}$ 上的 $n$ 阶矩阵，求证：$A$ 相似于一个 $\mathbb{K}$ 上主对角元全为零的矩阵的充要条件是$\mathrm{tr}(A)=0$.

\item\label{proposition:迹零阵必相似于主对角全为零的矩阵-2} 设 $A$ 是数域 $\mathbb{K}$ 上的 $n$ 阶矩阵，求证：$A$ 正交相似于一个 $\mathbb{R}$ 上主对角元全为零的矩阵的充要条件是$\mathrm{tr}(A)=0$.

\item\label{proposition:迹零阵必相似于主对角全为零的矩阵-3} 设 $C$ 是数域 $\mathbb{K}$ 上的 $n$ 阶矩阵，求证：存在 $\mathbb{K}$ 上的 $n$ 阶矩阵 $A,B$，使得 $AB - BA = C$ 的充要条件是 $\mathrm{tr}(C)=0$。

\item\label{proposition:迹零阵必相似于主对角全为零的矩阵-4} $\mathrm{dim}\left\{ AB-BA|A,B\in \mathbb{K} ^{n\times n} \right\} =\mathrm{dim}\left\{ A|\mathrm{tr}\left( A \right) =0,A\in \mathbb{K}^{n\times n} \right\} =n^2-1.$
\end{enumerate}
\end{proposition}
\begin{proof}
\begin{enumerate}[(1)]
\item {\heiti 必要性:}由矩阵的迹是矩阵的相似不变量,故$\mathrm{tr}(A)=0$.

{\heiti 充分性:}对阶数进行归纳。当 $n = 1$ 时，$A = O$，结论显然成立。设阶数小于 $n$ 时结论成立，现证 $n$ 阶的情形。由于题目的条件和结论在相似关系下不改变，故不妨从一开始就假设 $A$ 是有理标准型
\begin{align*}
\boldsymbol{F}=\mathrm{diag}\{\boldsymbol{F}(d_1(\lambda)),\cdots,\boldsymbol{F}(d_k(\lambda))\}
\end{align*}
其中 $d_i(\lambda)$ 是 $A$ 的非常数不变因子，$d_i(\lambda)\mid d_{i + 1}(\lambda) (1\leqslant  i\leqslant  k - 1)$，$\deg d_i(\lambda)=r_i$。若 $r_i$ 都为 $1$，则 $d_1(\lambda)=\cdots = d_n(\lambda)=\lambda - c$，从而 $A = c\boldsymbol{I}_n$。又 $\mathrm{tr}(A)=0$，故 $c = 0$，从而 $A = O$，结论成立。以下假设存在某个 $r_i>1$，将第 $(1,1)$ 分块与第 $(i,i)$ 分块对换，这是一个相似变换，此时矩阵的第 $(1,1)$ 元为零，故不妨设 $A$ 的第 $(1,1)$ 元为零。注意到矩阵 $A=\begin{pmatrix}
0 & \boldsymbol{\alpha}' \\
\boldsymbol{\beta} & \boldsymbol{B}
\end{pmatrix}$，其中 $\boldsymbol{\alpha},\boldsymbol{\beta}\in\mathbb{K}^{n - 1}$，$\boldsymbol{B}\in M_{n - 1}(\mathbb{K})$，$\mathrm{tr}(\boldsymbol{B})=0$。由归纳假设，存在 $\mathbb{K}$ 上的 $n - 1$ 阶非异阵 $\boldsymbol{Q}$，使得 $\boldsymbol{Q}^{-1}\boldsymbol{B}\boldsymbol{Q}$ 的主对角元全为零，令 $\boldsymbol{P}=\begin{pmatrix}
1 & \boldsymbol{O} \\
\boldsymbol{O} & \boldsymbol{Q}
\end{pmatrix}$ 为 $\mathbb{K}$ 上的 $n$ 阶非异阵，则 $\boldsymbol{P}^{-1}\boldsymbol{A}\boldsymbol{P}=\begin{pmatrix}
0 & \boldsymbol{\alpha}'\boldsymbol{Q} \\
\boldsymbol{Q}^{-1}\boldsymbol{\beta} & \boldsymbol{Q}^{-1}\boldsymbol{B}\boldsymbol{Q}
\end{pmatrix}$ 的主对角元全为零，结论得证。

\item {\heiti 必要性：}若矩阵 \( A \) 正交相似于主对角元全为零的矩阵，则由相似变换不改变迹可知, \(\mathrm{tr}(A) = 0\)。

{\heiti 充分性:}（i）若$A$有一个对角元为$0$，不妨设$A$的左上角元为$0$，则可设
\begin{align*}
A=\begin{pmatrix}
0& \beta^T\\
\alpha& A_1\\
\end{pmatrix},
\end{align*}
假设结论对$n-1$成立，则存在正交矩阵$P$，使得
$P^TA_1P=\Lambda_1$，其中$\Lambda_1$的对角元全为$0$。
于是
\begin{align*}
\begin{pmatrix}
1& 0\\
0& P^T\\
\end{pmatrix} A \begin{pmatrix}
1& 0\\
0& P\\
\end{pmatrix} =\begin{pmatrix}
0& *\\
*& \Lambda_1\\
\end{pmatrix}.
\end{align*}
此时$\begin{pmatrix}
1& 0\\
0& P^T\\
\end{pmatrix}$就是所求正交矩阵。


(ii)若$A$的对角元均不为$0$，则由$\text{tr}(A)=0$知，至少存在两个异号的对角元。
不妨设$A$的二阶右下角主子阵为
\begin{align*}
\begin{pmatrix}
a_{n-1}& x\\
y& a_n\\
\end{pmatrix},
\end{align*}
其中$a_{n-1}a_n<0$。
待定正交矩阵$Q=\begin{pmatrix}
I_{n-2}& & \\
& -b& a\\
& a& b\\
\end{pmatrix}$，其中$a^2+b^2=1$。则
\begin{align*}
Q^TAQ=\begin{pmatrix}
I_{n-2}& & \\
& -b& a\\
& a& b\\
\end{pmatrix}^T A \begin{pmatrix}
I_{n-2}& & \\
& -b& a\\
& a& b\\
\end{pmatrix}
\end{align*}
的二阶右下角主子阵为
\begin{align}
\begin{pmatrix}
-b& a\\
a& b\\
\end{pmatrix} \begin{pmatrix}
a_{n-1}& x\\
y& a_n\\
\end{pmatrix} \begin{pmatrix}
-b& a\\
a& b\\
\end{pmatrix} =\begin{pmatrix}
*& *\\
*& a^2a_{n-1}+ab(x+y)+b^2a_n\\
\end{pmatrix}.\label{eq::::::32r3::::--230fj23rt23t1}
\end{align}
记
\begin{align*}
f(\theta)=a_{n-1}\cos^2\theta+(x+y)\cos\theta\sin\theta+a_n\sin^2\theta,
\end{align*}
则$f(0)f\left( \frac{\pi}{2} \right) =a_{n-1}a_n<0$.由零点存在定理知，存在$\theta_0\in\left(0,\frac{\pi}{2}\right)$，使得$f(\theta_0)=0$。
于是取$a=\cos\theta_0$，$b=\sin\theta_0,$则
$$a^2a_{n-1}+ab(x+y)+b^2a_n=f(\theta_0)=0.$$
从而由\eqref{eq::::::32r3::::--230fj23rt23t1}式知
\begin{align*}
Q^TAQ=\begin{pmatrix}
*& *\\
*& 0\\
\end{pmatrix}.
\end{align*}
再由(i)可知结论成立。

\item 必要性由矩阵迹的线性和交换性即得，下证充分性。

由于题目的条件和结论在同时相似变换 $A\mapsto P^{-1}AP$，$B\mapsto P^{-1}BP$，$C\mapsto P^{-1}CP$ 下不改变，故由\nrefpro{proposition:迹零阵必相似于主对角全为零的矩阵}{(1)}不妨从一开始就假设 $C=(c_{ij})$ 的主对角元 $c_{ii}=0 (1\leqslant  i\leqslant  n)$。取定 $A = \mathrm{diag}\{\lambda_1,\lambda_2,\cdots,\lambda_n\}$ 为 $\mathbb{K}$ 上的主对角元互异的对角矩阵。设 $B=(x_{ij})$，则 $AB - BA = C$ 等价于方程 $\lambda_ix_{ij}-\lambda_jx_{ij}=c_{ij}$。当 $i = j$ 时，上式恒成立，故 $x_{ii}$ 可任取。当 $i\neq j$ 时，$x_{ij}=\frac{c_{ij}}{\lambda_i - \lambda_j}$ 被唯一确定。因此，一定存在 $\mathbb{K}$ 上的矩阵 $A,B$，使得 $AB - BA = C$ 成立.

\item 由\ref{proposition:迹零阵必相似于主对角全为零的矩阵-3}可知$\mathrm{dim}\left\{ AB-BA|A,B\in \mathbb{K} ^{n\times n} \right\} =\mathrm{dim}\left\{ A|\mathrm{tr}\left( A \right) =0,A\in \mathbb{K}^{n\times n} \right\} $,故只须证
\begin{align*}
\mathrm{dim}\left\{ A|\mathrm{tr}\left( A \right) =0,A\in \mathbb{K}^{n\times n} \right\} =n^2-1.
\end{align*}
{\color{blue} 证法一:}
注意到$\{ A|\mathrm{tr}(A) =0,A\in \mathbb{K}^{n\times n} \}$的一组基为  
\begin{align*}
E_{ij}(i\ne j),\quad E_{ii}-E_{i+1,i+1}(i=1,2,\cdots,n-1),
\end{align*}
并且这组基中共含有$n^2-1$个线性无关的矩阵．故$\mathrm{dim}\left\{ A|\mathrm{tr}\left( A \right) =0,A\in \mathbb{K}^{n\times n} \right\} =n^2-1.$  

{\color{blue} 证法二:}将矩阵的迹看作$\mathbb{K} ^{n\times n}\rightarrow \mathbb{K}$的一个线性映射,记作$\mathrm{tr}$．显然$\mathrm{tr}$是满射,于是  
\begin{align*}
\mathrm{dim}\{ A|\mathrm{tr}(A) =0,A\in \mathbb{K}^{n\times n} \} = \mathrm{dim}\,\mathrm{Ker}\,\mathrm{tr} = \mathrm{dim}\mathbb{K} ^{n\times n} - \mathrm{dim}\,\mathrm{Im}\,\mathrm{tr} = n^2 - \mathrm{dim}\,\mathbb{K} = n^2-1.
\end{align*}
\end{enumerate}
\end{proof}

\begin{proposition}
设$\mathrm{tr}(A)=0$，$f(x)$是多项式且$\deg f(x)\geqslant  1$，证明：存在可逆方阵$X,Y$使得$A = f(XY)-f(YX)$.
\end{proposition}
\begin{note}
另外的证明见清疏讲义.
\end{note}
\begin{proof}
由\subcref{proposition:迹零阵必相似于主对角全为零的矩阵}{proposition:迹零阵必相似于主对角全为零的矩阵-4}可知，存在可逆矩阵$P$，使得
\begin{align*}
A = P\Lambda P^{-1},
\end{align*}
其中$\Lambda = (b_{ij})_{n\times n}$且$b_{ii} = 0\ (i = 1,2,\cdots,n)$。注意到条件与结论在线性变换$A \mapsto P\Lambda P^{-1}$下不改变，因此不妨设$A = \Lambda$，则取$Y = \begin{pmatrix}
1 & & & \\
& 2 & & \\
& & \ddots & \\
& & & n
\end{pmatrix}$，设$X = (x_{ij})$，记矩阵$U = XY$，$A_1 = AY$，则$A_1$也是主对角元全为零的矩阵，设$A_1 = (a_{ij})_{n\times n}$，其中$a_{ii} = 0\ (i = 1,2,\cdots,n)$。从而
\begin{align}
&A = f(XY) - f(YX) \Longleftrightarrow A = f(U) - f(YUY^{-1}) = f(U) - Yf(U)Y^{-1} \notag \\
&\Longleftrightarrow AY = f(U)Y - Yf(U) \Longleftrightarrow A_1 = f(U)Y - Yf(U). \label{100.49}
\end{align}
由\subcref{proposition:迹零阵必相似于主对角全为零的矩阵}{proposition:迹零阵必相似于主对角全为零的矩阵-4}可知，存在$M_0 \in \mathbb{C}^{n\times n}$，使得$A_1 = M_0Y - YM_0$。令$M_1 = M_0 + \lambda Y$，$\forall \lambda \in \mathbb{C}$，则
\begin{align}
A_1 = M_1Y - YM_1. \label{100.48}
\end{align}
设$M_0 = (x_{ij})_{n\times n}$，则$M_1 = \begin{pmatrix}
x_{11} + \lambda & x_{12} & \cdots & x_{1n} \\
x_{21} & x_{22} + 2\lambda & \cdots & x_{2n} \\
\vdots & \vdots & \ddots & \vdots \\
x_{n1} & x_{n2} & \cdots & x_{nn} + n\lambda
\end{pmatrix}$。从而可设$D_k: |z - x_{kk} - k\lambda| \leqslant \sum_{j\ne k} x_{kj}\ (1 \leqslant k \leqslant n)$为$M_1$的$n$个戈氏圆盘。取$\lambda$充分大，使得
\begin{align}
|x_{kk} - k\lambda - f(0)| > \sum_{j\ne k} x_{kj}\ (1 \leqslant k \leqslant n), \label{100.46}
\end{align}
并且$D_1, \cdots, D_n$互不相交，从而此时$D_i\ (1 \leqslant i \leqslant n)$各自构成一个连通分支。于是由第二圆盘定理可知，每个连通分支$D_i$中有且仅有一个特征值，进而$M_1$此时有$n$个不同的特征值，故此时$M_1$可对角化。于是存在可逆阵$G$，使得
\begin{align*}
M_1 = G\begin{pmatrix}
\lambda_1 & & & \\
& \lambda_2 & & \\
& & \ddots & \\
& & & \lambda_n
\end{pmatrix}G^{-1},
\end{align*}
其中$\lambda_1, \cdots, \lambda_n$为$M_1$的所有特征值。取$u_i \in \mathbb{C}$为$f(x) = \lambda_i\ (1 \leqslant i \leqslant n)$的根，又由\eqref{100.46}式可知，$f(0)$落在$\bigcup_{i=1}^n D_i$之外，故$M_1$的所有特征值$\lambda_1, \cdots, \lambda_n \ne f(0)$。从而$u_i \ne 0\ (1 \leqslant i \leqslant n)$。于是令$U = G\begin{pmatrix}
u_1 & & & \\
& u_2 & & \\
& & \ddots & \\
& & & u_n
\end{pmatrix}G^{-1}$，则
\begin{align*}
f(U) = G\begin{pmatrix}
f(u_1) & & & \\
& f(u_2) & & \\
& & \ddots & \\
& & & f(u_n)
\end{pmatrix}G^{-1} = G\begin{pmatrix}
\lambda_1 & & & \\
& \lambda_2 & & \\
& & \ddots & \\
& & & \lambda_n
\end{pmatrix}G^{-1} = M_1.
\end{align*}
再结合\eqref{100.48}式可知$A_1 = f(U)Y - Yf(U)$。令$X = UY^{-1}$，则再根据\eqref{100.49}式就可知
\begin{align*}
A = f(XY) - f(YX).
\end{align*}
故结论得证。
\end{proof}

\begin{example}
设\(A,B\)是\(n\)阶复方阵,且\(A\)可逆. 已知存在正整数\(t\),使得对\(k \in \{t+1, t+2, \cdots, t+n\}\),\(B\)与\(A^k\)均可交换. 证明: \(B\)与\(A\)可交换.
\end{example}
\begin{proof}
设\(A\)的特征多项式为\(p(\lambda) = \lambda^n + c_{n-1}\lambda^{n-1} + \cdots + c_1\lambda + c_0\). 因为\(A\)可逆,所以\(c_0 \neq 0\). 由Cayley-Hamilton定理,\(p(A) = 0\),即
\begin{align*}
A^n + c_{n-1}A^{n-1} + \cdots + c_1A + c_0I = 0.
\end{align*}
将上式两边同乘\(A^t\),得到
\begin{align*}
A^{n+t} + c_{n-1}A^{n+t-1} + \cdots + c_1A^{t+1} + c_0A^t = 0.
\end{align*}
移项可得
\begin{align*}
A^t = -\frac{1}{c_0}(A^{n+t} + c_{n-1}A^{n+t-1} + \cdots + c_1A^{t+1}).
\end{align*}
这说明\(A^t\)是\(A^{t+1}, A^{t+2}, \cdots, A^{t+n}\)的线性组合. 因为\(B\)与\(A^{t+1}, A^{t+2}, \cdots, A^{t+n}\)均可交换,所以\(B\)与\(A^t\)也可交换,即
\begin{align*}
BA^t = A^tB.
\end{align*}
又因为\(B\)与\(A^{t+1}\)可交换,即\(BA^{t+1} = A^{t+1}B\). 将左边展开得
\begin{align*}
BA^{t+1} = BA^tA = A^tBA.
\end{align*}
因此\(A^tBA = A^{t+1}B\). 由于\(A^tBA = BA^tA\),故\(A^t(BA - AB) = 0\). 因为\(A\)可逆,所以\(A^t\)可逆,左乘\(A^{-t}\)可得\(BA - AB = 0\),即
\begin{align*}
BA = AB.
\end{align*}
所以\(B\)与\(A\)可交换.
\end{proof}








\end{document}